% CLAUDE REPORT::B3-RELATIVIZATION-0001::2026-11-10_2151_PST
\documentclass[11pt]{article}
\usepackage[margin=0.9in]{geometry}
\usepackage{amsmath,amssymb,amsthm,hyperref}
\newtheorem{theorem}{Theorem}
\newtheorem{lemma}[theorem]{Lemma}
\newtheorem{corollary}[theorem]{Corollary}
\newtheorem{proposition}[theorem]{Proposition}
\theoremstyle{definition}
\newtheorem{definition}[theorem]{Definition}
\newtheorem{remark}[theorem]{Remark}
\title{B3-RELATIVIZATION-0001: Where the Conditional Bridge Stands,\\
and Why No Relativizing or Algebrizing Argument Can Prove It}
\author{Claude research agent (additive; not an LSPV role)}
\date{November 10, 2026, 21:51 PST}
\begin{document}
\maketitle

\noindent\textbf{Status summary.}
Theorems~\ref{thm:dnf}, \ref{thm:mono} and~\ref{thm:rel} are \textbf{PROVED/EXACT}
(they are short and follow from standard definitions).
Corollary~\ref{cor:known} is \textbf{PROVED/EXACT} given cited \textbf{KNOWN THEOREMS}.
B3 for the project system $R_{\mathrm{mix341}}$, the superpolynomial lower bound on
$\overline W_{R_{\mathrm{mix341}}}$, and P versus NP all remain \textbf{OPEN}.
This note is additive. It changes no schedules, roles, ownership regions or allocations,
and it does not edit root \texttt{main.tex}, because this agent owns no region there.

\section{Setting: the shared target}
Fix a Cook--Reckhow proof system $R$ for CNF unsatisfiability. That means $R$ is
sound and complete, and a polynomial-time checker verifies its proofs; we write
$\mathrm{cost}_R(\pi)\ge 1$ per derived node, as in the 2026-11-02 directive. Define
\[
W_R(F)=\min_{\pi:F\vdash_R\bot}\mathrm{cost}_R(\pi),\qquad
\overline W_R(N)=\max_{F\in\mathrm{UNSAT},\ |\langle F\rangle|\le N}W_R(F).
\]
The inner minimum lets the best derivation be chosen for free. As the user's directive
says, this is a legitimate lower-bound strategy: discovering the derivation is never
charged, so any lower bound on $W_R$ bounds every procedure that \emph{executes} an
$R$-derivation. Write
\[
\mathrm{PB}_R:\ \exists C,k\ \forall N\ \overline W_R(N)\le C(N+1)^k
\qquad\text{($R$ is polynomially bounded in node cost)},
\]
\[
\mathrm{B3}_R:\ \mathrm{SAT}\in\mathrm P\ \Longrightarrow\ \mathrm{PB}_R .
\]
Assume node cost and proof bit length are polynomially related, up to the encoding of a
node. This holds for every system in the corpus, because each node has a
polynomial-size description and the checker is polynomial-time. Then $\mathrm{PB}_R$ is
the usual statement that $R$ is p-bounded, and $\mathrm{PB}_R\Rightarrow\mathrm{NP}=\mathrm{coNP}$
by Cook--Reckhow~\cite{CR79}.

\section{Results}

\begin{theorem}[Disjunctive normal form of B3]\label{thm:dnf}
For every $R$, $\mathrm{B3}_R\iff(\mathrm P\ne\mathrm{NP})\vee\mathrm{PB}_R$.
\end{theorem}
\begin{proof}
This is propositional logic: $(a\Rightarrow b)\equiv(\neg a\vee b)$, with $a$ being
$\mathrm{SAT}\in\mathrm P$, which is equivalent to $\mathrm P=\mathrm{NP}$ by Cook--Levin.
\end{proof}

\begin{corollary}[With a proved lower bound, B3 \emph{is} P $\ne$ NP]\label{cor:known}
Suppose $\neg\mathrm{PB}_R$ is a theorem, meaning a superpolynomial lower bound on
$\overline W_R$ has been proved. Then $\mathrm{B3}_R$ is provably equivalent to
$\mathrm P\ne\mathrm{NP}$. In particular:
\begin{enumerate}
\item For ordinary Resolution, $\mathrm{B3}_{\mathrm{Res}}\iff\mathrm P\ne\mathrm{NP}$,
because Haken's pigeonhole lower bound~\cite{Haken85} is a KNOWN THEOREM, transferred to
3-CNF by standard local encoding.
\item The same holds for Cutting Planes, Polynomial Calculus, and bounded-depth Frege,
since each has a KNOWN superpolynomial lower bound. If Braun's 2026 dag-like
Res$(\oplus)$ lower bound for bit-PHP~\cite{Braun26} survives review, it holds for
Res$(\oplus)$ too. That bound is recorded in the corpus at LSPV-1/369 as a preprint
pending review.
\end{enumerate}
So no bridge B3 is ``weaker than P $\ne$ NP'' for any system whose lower bound is
already known. For such $R$, proving B3 is proving the separation itself.
\end{corollary}

\begin{theorem}[Monotonicity: B3 gets easier exactly as the lower bound gets harder]\label{thm:mono}
Suppose $R'$ p-simulates $R$ in node cost: there is a polynomial $q$ with
$W_{R'}(F)\le q(|\langle F\rangle|+W_R(F))$ for all UNSAT $F$. Then
$\mathrm{PB}_R\Rightarrow\mathrm{PB}_{R'}$. Hence
$\mathrm{B3}_R\Rightarrow\mathrm{B3}_{R'}$ and $\neg\mathrm{PB}_{R'}\Rightarrow\neg\mathrm{PB}_R$.
\end{theorem}
\begin{proof}
$\overline W_{R'}(N)\le q(N+\overline W_R(N))$. For B3, combine this with
Theorem~\ref{thm:dnf}.
\end{proof}

\noindent The pair (lower bound for $R$, $\mathrm{B3}_R$) therefore trades difficulty
monotonically along the simulation order. Strengthening $R$ (for example, adding
VARDEF/extension, which $R_{\mathrm{mix341}}$ already has) moves work from the bridge to
the lower bound. The lower bound for an $R$ that p-simulates Extended Resolution is the
long-open EF lower-bound problem. No choice of $R$ makes both halves easy unless
P $\ne$ NP is already known.

\begin{definition}[Relativized bridge]
For an oracle $O$, let $\mathrm{B3}^O_R$ be the statement
$\mathrm{SAT}\in\mathrm P^O\Rightarrow\mathrm{PB}_R$. Equivalently, take
$\mathrm{NP}^O\subseteq\mathrm P^O$ as the hypothesis; the proof below works for both
readings. $R$ stays the same fixed, oracle-free proof system. A proof of $\mathrm{B3}_R$
\emph{relativizes} if the same argument proves $\mathrm{B3}^O_R$ for every oracle $O$.
\end{definition}

\begin{theorem}[Relativization and algebrization firewall for B3]\label{thm:rel}
For every fixed Cook--Reckhow system $R$, the following are equivalent:
\begin{enumerate}
\item[(i)] $\mathrm{B3}^O_R$ holds for every oracle $O$;
\item[(ii)] $\mathrm{B3}^{O}_R$ holds for $O=\mathrm{TQBF}$;
\item[(iii)] $\mathrm{B3}^{\widetilde O}_R$ holds for $\widetilde O$, the multilinear
extension of TQBF over $\mathbb F_2$-extension fields, as in Aaronson--Wigderson~\cite{AW09};
\item[(iv)] $\mathrm{PB}_R$.
\end{enumerate}
Consequently, a relativizing or algebrizing proof of $\mathrm{B3}_R$, for \emph{any}
fixed sound and complete polynomially checkable $R$, would prove that $R$ is p-bounded.
That would give NP $=$ coNP. Unless NP $=$ coNP, every proof of B3 must be
non-relativizing and non-algebrizing.
\end{theorem}
\begin{proof}
(iv)$\Rightarrow$(i), (ii) and (iii): the conclusion holds outright.
(i)$\Rightarrow$(ii) is specialization.
(ii)$\Rightarrow$(iv): $\mathrm{SAT}\le_p\mathrm{TQBF}$, so $\mathrm{SAT}\in\mathrm P^{\mathrm{TQBF}}$.
Indeed $\mathrm{NP}^{\mathrm{TQBF}}=\mathrm{PSPACE}=\mathrm P^{\mathrm{TQBF}}$~\cite{BGS75}.
The hypothesis of $\mathrm{B3}^{\mathrm{TQBF}}_R$ is therefore true. Since $\overline W_R$
does not mention the oracle, its conclusion $\mathrm{PB}_R$ holds.
(iii)$\Rightarrow$(iv): $\widetilde O$ agrees with TQBF on Boolean points, so
$\mathrm{TQBF}\in\mathrm P^{\widetilde O}$. Also $\widetilde O$ is computable in PSPACE,
so $\mathrm P^{\widetilde O}=\mathrm{NP}^{\widetilde O}=\mathrm{PSPACE}$~\cite{AW09}.
The hypothesis is again true, and the oracle-free conclusion follows.
The final sentence follows from Cook--Reckhow~\cite{CR79}.
\end{proof}

\begin{remark}[What the firewall does and does not say]
\begin{itemize}
\item It is not the Baker--Gill--Solovay barrier restated. BGS concerns a proof of
P $\ne$ NP. Theorem~\ref{thm:rel} applies to $\mathrm{B3}_R$ even for systems such as
$R_{\mathrm{mix341}}$ and EF, where no lower bound is known and where B3 is not known to
be equivalent to P $\ne$ NP.
\item It rules out every ``black-box'' computation-to-proof transfer. Such a transfer
uses only the input/output behaviour or running time of a SAT decider and would work
equally well for an oracle machine. The corpus's answer-only transducer firewall
(LSPV-2/478) and Gödel wrapper firewall (LSPV-1/357) are special cases of this pattern.
\item It is consistent with the known white-box sufficient route. Cook's translation and
Krajíček's theorem~\cite{Cook75,Kra95} say that if PV proves that a specific algorithm
$A$ solves SAT, then $A$ yields polynomial EF refutations uniformly. That route reads
$A$'s \emph{code and a feasible correctness proof}, and it does not relativize to
$\mathrm{TQBF}$-oracle machines with oracle-free EF. So any successful B3 argument must
use the internal structure of the hypothesised algorithm, such as its transition
function, its locality, or a feasible correctness proof.
\item If one also lets $R$'s checker query the oracle ($R^O$), then
$R^{\mathrm{TQBF}}$ is trivially p-bounded and the theorem says nothing. The firewall is
specific to a \emph{preselected oracle-free} $R$, which is exactly the directive's
requirement.
\end{itemize}
\end{remark}

\section{Consequences for the multisign DAG-hypergraph programme}
The user's method is sound for its own target. It builds the signed DAG hypergraph of a
3-CNF, defines the minimal reduced logical state, and minimises work over all valid
reduction sequences, with the best sequence granted for free. That yields a lower bound
on $W_R(F)$ for the fixed reduction calculus $R$, and on $\overline W_R(N)$ in the worst
case. The present results locate the remaining difficulty precisely:
\begin{enumerate}
\item A superpolynomial $\overline W_R$ bound \emph{plus} $\mathrm{B3}_R$ is a proof of
P $\ne$ NP (Theorem~\ref{thm:dnf}). The free-sequence minimum adds no hidden cost, but it
also cannot make the bridge cheaper.
\item If the reduction calculus is weak enough that hard families are already known (as
for Resolution-like reductions: Tseitin expanders, PHP), the bridge is the entire
problem (Corollary~\ref{cor:known}). If it is strong enough to include extension
variables, the lower bound is the open EF problem (Theorem~\ref{thm:mono}).
\item Whatever the choice, the bridge must be proved by a white-box argument about the
hypothesised polynomial SAT algorithm (Theorem~\ref{thm:rel}). Any argument that treats
the solver only through its answers or running time is ruled out in advance.
\end{enumerate}

\section*{Falsifiers checked}
\begin{itemize}
\item[F1] \emph{Is there an oracle making SAT easy while $\overline W_R$ stays
superpolynomial?} Yes: TQBF with $R=$ Resolution, using BGS and Haken (both KNOWN). This
confirms that (ii)$\Rightarrow$(iv) is not vacuous for weak $R$.
\item[F2] \emph{Could $\overline W_R$ secretly depend on the oracle?} No. $R$'s checker
is a fixed oracle-free polynomial-time machine, so the set of legal proofs is unchanged.
\item[F3] \emph{Does the algebrizing case need extra care?} Yes, and it is handled:
$\widetilde O\in\mathrm{PSPACE}$ gives $\mathrm{NP}^{\widetilde O}\subseteq\mathrm{PSPACE}$,
and TQBF$\le\widetilde O$ gives $\mathrm{PSPACE}\subseteq\mathrm P^{\widetilde O}$.
\end{itemize}

\section*{Corpus reviewed}
\texttt{AGENTS.md}; the five-role worst-case directive
\texttt{research/coordination/2026-11-02\_1204\_WORST\_CASE\_MINIMUM\_DERIVATION\_FIVE\_ROLE\_REFOCUS.md};
LSPV-1/357 (Gödel reflection B3 firewall); LSPV-1/369 (Braun Res$(\oplus)$ transfer to
the affine fragment); the \texttt{proof2.txt} node [P2-B2] (noncircular B3, OPEN) and
LSPV-2/478 (answer-only oracle firewall, as recorded in the LSPV-5/479 proof2
pre-write); the newest commit on \texttt{main} (LSPV-3/561).

\section*{Duplication / contradiction audit}
Repository search for a relativized B3, a TQBF oracle with B3, or ``B3 equivalent to
P $\ne$ NP'' found no prior statement of Theorems~\ref{thm:dnf}--\ref{thm:rel}.
LSPV-1/357 proves a different, single-wrapper incompleteness firewall and gives the
PV/EF sufficient condition, which is reused here. LSPV-2/478 is a resource-bounded
special case of the black-box pattern. No contradiction with any verified result was
found.

\section*{What changed relative to root main}
Nothing in \texttt{main.tex} is edited. Proposed evidence for [P2-B2]: B3 for any
non-p-bounded fixed $R$ requires a non-relativizing, non-algebrizing, white-box argument.
For any $R$ with a proved lower bound, B3 is equivalent to P $\ne$ NP. LSPV-5 may adopt
or reject this at its next reconciliation.

\section*{Adopted / rejected / retired prior components}
Adopted: the $\overline W_R$ definition and the B1/B2/B3 hierarchy; LSPV-1/357's PV/EF
route.
Retired as routes to B3, though not as diagnostics: every answer-only or runtime-only
transfer.

\section*{Stale or overstated claims found}
None found. The directive's label ``weakest sufficient bridge'' is accurate within the
B1/B2/B3 order. Corollary~\ref{cor:known} adds that, for already-hard systems, it is not
weaker than the separation itself.

\section*{Literature queries (2026-11-10)}
``P = NP implies Extended Frege polynomially bounded bounded arithmetic PV'': this
surfaced Krajíček's characterisation (PV proves P $=$ NP iff PV defines an algorithm
returning a satisfying assignment or an EF refutation) and Pich--Santhanam-style
consistency results (arXiv:1905.12935). ``relativized proof complexity oracle P=NP'':
this surfaced the oracle constructions for optimal proof systems (arXiv:2408.07408) and
relativized optimal proof systems (arXiv:2602.02294). None of these states
Theorem~\ref{thm:rel}.

\begin{thebibliography}{9}
\bibitem{AW09} S.~Aaronson, A.~Wigderson, Algebrization: A New Barrier in Complexity Theory, ACM TOCT 1(1), 2009.
\bibitem{BGS75} T.~Baker, J.~Gill, R.~Solovay, Relativizations of the P =? NP question, SIAM J.\ Comput.\ 4 (1975).
\bibitem{Braun26} K.~Braun, An Exponential Lower Bound for the Bit Pigeonhole Principle in Resolution over Parities, arXiv:2609.23015 (2026), preprint.
\bibitem{Cook75} S.~A.~Cook, Feasibly constructive proofs and the propositional calculus, STOC 1975.
\bibitem{CR79} S.~A.~Cook, R.~A.~Reckhow, The relative efficiency of propositional proof systems, JSL 44 (1979).
\bibitem{Haken85} A.~Haken, The intractability of resolution, TCS 39 (1985).
\bibitem{Kra95} J.~Krajíček, Bounded Arithmetic, Propositional Logic and Complexity Theory, CUP 1995.
\end{thebibliography}

\noindent P versus NP remains OPEN.
\end{document}
% END_OF_VERIFIED_REPORT::CLAUDE::B3-RELATIVIZATION-0001
